Gravitational radiation in empty asymptotically flat 
space has been examined by means of asymptotic expansions 
by a number of authors.\,(1--4) They find that the different 
components of the outgoing radiation field ``peel off'', that 
is, they go as different powers of the affine radial distance.  
If one wishes to investigate how this behaviour is modified 
by the presence of matter, one is faced with a difficulty 
that does not arise in the case of, say, electromagnetic 
radiation in matter. For this one can consider the radiation 
travelling through an infinite uniform medium that is static 
apart from the disturbance created by the radiation. In the 
case of gravitational radiation this is not possible. For, 
if the medium were initially static, its own self–gravitation 
would cause it to contract in on itself and it would cease to 
be static. Hence one is forced to investigate gravitational 
radiation in matter that is either contracting or expanding.  

As in Chapter 2, we identify the Weyl or conformal 
tensor $C_{abcd}$ with the free gravitational field and the 
Ricci–tensor $R_{ab}$ with the contribution of the matter to the 
curvature. Instead of considering gravitational radiation in asymptotically flat space, that is, space that approaches 
flat space at large radial distances, we consider it in 
asymptotically conformally flat space. As it is only 
conformally flat, the Ricci–tensor and the density of matter 
need not be zero.  

To avoid essentially non–gravitational phenomena such 
as sound waves, we will consider gravitational radiation 
travelling through dust. It was shown in Chapter 2 that a 
conformally flat universe filled with dust must have one of 
the metrics:  

(a) 
\[
ds^{2} = -\Omega^{2}\left(d\iota^{2} - d\varrho^{2} - \sin^{2}\varrho \left(d\theta^{2} + \sin^{2}\theta\, d\phi^{2}\right)\right),
\qquad \Omega = A(1 - \cos \tau) \tag{1.1}
\]

(b) 
\[
ds^{2} = \Omega^{2}\left(dt^{2} - d\varrho^{2} - \varrho^{2}\left(d\theta^{2} + \sin^{2}\theta\, d\phi^{2}\right)\right),
\qquad \Omega = \tfrac{1}{6}A t^{2} \tag{1.2}
\]

(c) 
\[
ds^{2} = \Omega^{2}\left(dt^{2} - d\varrho^{2} - \sinh^{2}\varrho\left(d\theta^{2} + \sin^{2}\theta\, d\phi^{2}\right)\right),
\qquad \Omega = A(\cosh t - 1) \tag{1.3}
\]

Type (a) represents a universe in which the matter 
expands from the initial singularity with insufficient energy 
to reach infinity and so falls back again to another 
singularity. It is therefore unsuitable for a discussion of gravitational radiation by a method of asymptotic expansions 
since one cannot get an infinite distance from the source.  

Type (b) is the Einstein–de Sitter universe in which the 
matter has just sufficient energy to reach infinity. It is 
thus a special case. D. Norman\,(5) has investigated the 
``peeling off'' behaviour in this case using Penrose’s conformal 
technique\,(6). He was however forced to make certain assumptions 
about the movement of the matter which will be shown to 
be false. Moreover, he was misled by the special nature of 
the Einstein–de Sitter universe in which affine and luminosity 
distances differ. Another reason for not considering radiation 
in the Einstein–de Sitter universe is that it is unstable.  
The passage of a gravitational wave will cause it to contract 
again eventually and develop a singularity.  

We will therefore consider radiation in a universe of 
type (c) which corresponds to the general case where the 
matter is expanding with more than enough energy to avoid 
contracting again.  

\sectiontoc{2. The Newman–Penrose Formalism}

We employ the notation of Newman and Penrose.\,(3) A 
tetrad of null vectors, $l^{\mu}, n^{\mu}, m^{\mu}, \bar{m}^{\mu}$ is introduced
where:
\[
l_{\mu}l^{\mu} = n_{\mu}n^{\mu} = m_{\mu}m^{\mu} = \bar{m}_{\mu}\bar{m}^{\mu} =
l_{\mu}m^{\mu} = l_{\mu}\bar{m}^{\mu} = n_{\mu}m^{\mu} = n_{\mu}\bar{m}^{\mu} = 0,
\]
\[
l_{\mu}n^{\mu} = 1, \qquad m_{\mu}\bar{m}^{\mu} = -1.
\]

We label these vectors with a tetrad index
\[
Z^{\mu}_{a} = (l^{\mu}, n^{\mu}, m^{\mu}, \bar{m}^{\mu}), \qquad a = 1,2,3,4.
\]

Tetrad indices are raised and lowered with the metric
\[
\eta_{ab} =
\begin{bmatrix}
0 & 1 & 0 & 0 \\
1 & 0 & 0 & 0 \\
0 & 0 & 0 & -1 \\
0 & 0 & -1 & 0
\end{bmatrix}. \tag{2.1}
\]

We have,
\[
g^{\mu\nu} = \eta^{ab} Z^{\mu}_{a} Z^{\nu}_{b}
= l^{\mu}n^{\nu} + n^{\mu}l^{\nu} - m^{\mu}\bar{m}^{\nu} - \bar{m}^{\mu}m^{\nu}. \tag{2.2}
\]

Ricci rotation coefficients are defined by:
\[
\gamma_{abc} = Z^{\mu}_{a;\nu} Z^{\nu}_{b} Z_{c\mu}, \tag{2.3}
\]
\[
\gamma_{abc} = -\gamma_{acb}.
\]
In fact it is more convenient to work in terms of twelve 
complex combinations of rotation coefficients defined as follows:

\[
\kappa = \gamma_{131} = l_{\mu;\nu} m^{\mu} l^{\nu},
\]
\[
\pi = -\gamma_{241} = -n_{\mu;\nu} \bar{m}^{\mu} l^{\nu},
\]
\[
\epsilon = \tfrac{1}{2}(\gamma_{121} + \gamma_{341})
= \tfrac{1}{2}(l_{\mu;\nu} n^{\mu} l^{\nu} - m_{\mu;\nu} \bar{m}^{\mu} l^{\nu}),
\]
\[
\rho = \gamma_{134} = l_{\mu;\nu} m^{\mu} \bar{m}^{\nu},
\]
\[
\lambda = -\gamma_{244} = -n_{\mu;\nu} \bar{m}^{\mu} \bar{m}^{\nu},
\]
\[
\alpha = \tfrac{1}{2}(\gamma_{124} - \gamma_{344})
= \tfrac{1}{2}(l_{\mu;\nu} n^{\mu} \bar{m}^{\nu} - m_{\mu;\nu} \bar{m}^{\mu} \bar{m}^{\nu}),
\]
\[
\beta = \tfrac{1}{2}(\gamma_{123} - \gamma_{343})
= \tfrac{1}{2}(l_{\mu;\nu} n^{\mu} m^{\nu} - m_{\mu;\nu} \bar{m}^{\mu} m^{\nu}),
\]
\[
\sigma = \gamma_{133} = l_{\mu;\nu} m^{\mu} m^{\nu},
\]
\[
\mu = -\gamma_{243} = -n_{\mu;\nu} \bar{m}^{\mu} m^{\nu},
\]
\[
\nu = -\gamma_{242} = -n_{\mu;\nu} n^{\mu} \bar{m}^{\nu},
\]
\[
\gamma = \tfrac{1}{2}(\gamma_{122} - \gamma_{342})
= \tfrac{1}{2}(l_{\mu;\nu} n^{\mu} n^{\nu} - m_{\mu;\nu} \bar{m}^{\mu} n^{\nu}),
\]
\[
\tau = \gamma_{132} = l_{\mu;\nu} m^{\mu} n^{\nu}. \tag{2.4}
\]
\sectiontoc{3. Coordinates}

Like Newman and Penrose, we introduce a null coordinate 
\[
u (= X^{1}), \qquad g^{\mu\nu} u_{,\mu} u_{,\nu} = 0. \tag{3.1}
\]

We take 
\[
l_{\mu} = u_{,\mu}.
\]
Thus $l_{\mu}$ will be geodesic and irrotational. This implies
\[
\kappa = 0, \quad \rho = \bar{\rho}, \quad 
\epsilon = \bar{\epsilon}, \quad \alpha + \bar{\beta} = 0. \tag{3.2}
\]

We take $n^{\mu}, m^{\mu}, \bar{m}^{\mu}$ to be parallely transported 
along $l^{\mu}$. This gives
\[
\pi = \epsilon = 0. \tag{3.3}
\]

As a second coordinate we take an affine parameter $\gamma(=X^{2})$ 
along the geodesics $l^{\mu}$,
\[
\gamma_{,\mu} l^{\mu} = 1. \tag{3.4}
\]

$X^{3}$ and $X^{4}$ are two coordinates that label the geodesic 
in the surface $u=\text{const}$.  

\[
g^{\mu\nu} X^{i}{}_{,\mu} u_{,\nu} = X^{i}{}_{,\mu} l^{\mu} = 0. \tag{3.5}
\]

Thus,
\[
g^{\mu\nu} =
\begin{bmatrix}
0 & 1 & 0 & 0 \\
1 & 0 & g^{23} & g^{24} \\
0 & g^{23} & g^{33} & g^{34} \\
0 & g^{24} & g^{34} & g^{44}
\end{bmatrix}. \tag{3.6}
\]

In these coordinates
\[
l_{\mu} = \delta^{1}{}_{\mu}, \qquad l^{\mu} = \delta^{\mu}{}_{2},
\]
since $l_{\mu} l^{\mu} = 0$, $l_{\mu} m^{\mu} = 0$,  

\[
m^{\mu} = \cos\theta\, \delta^{\mu}{}_{3} + \sin\theta\, \delta^{\mu}{}_{4},
\]

\[
n^{\mu} = \delta^{\mu}{}_{1} + U \delta^{\mu}{}_{2} + X^{i} \delta^{\mu}{}_{i}, 
\qquad i=3,4. \tag{3.7}
\]
\sectiontoc{The Field Equations}

We may calculate the Ricci and Weyl tensor components from 
the relations
\[
R_{abcd} - R_{abdc} = \gamma_{ab[c;d]} - \gamma_{ab[d;c]}
+ \gamma_{aeb}\gamma^{e}{}_{cd} - \gamma_{aed}\gamma^{e}{}_{bc}, \tag{3.8}
\]

\[
R_{abcd} = C_{abcd} + \eta_{a[c}R_{d]b} + \eta_{b[d}R_{c]a}
+ \tfrac{R}{3}\eta_{a[c}\eta_{d]b}. \tag{3.9}
\]

Using the combinations of rotation coefficients already 
defined and with $\kappa = \pi = \epsilon = 0$, we have

\[
D\rho = \rho^{2} + \sigma \bar{\sigma} + \Phi_{00}, \tag{3.10}
\]
\[
D\sigma = 2\rho\sigma + \Psi_{0}, \tag{3.11}
\]
\[
D\tau = \rho\tau + \bar{\sigma}\pi + \Psi_{1} + \Phi_{01}, \tag{3.12}
\]
\[
D\alpha = \rho\alpha + \beta\bar{\sigma} + \Phi_{10}, \tag{3.13}
\]
\[
D\beta = \rho\beta + \alpha\sigma + \Psi_{1}, \tag{3.14}
\]
\[
D\gamma = \tau\alpha + \bar{\tau}\beta - \Psi_{2} - \Lambda + \Phi_{11}, \tag{3.15}
\]
\[
D\lambda = \rho\lambda + \mu\bar{\sigma} + \Phi_{20}, \tag{3.16}
\]
\[
D\mu = \rho\mu + \sigma\lambda + \Psi_{2} + 2\Lambda, \tag{3.17}
\]
\[
D\nu = \tau\lambda + \bar{\tau}\mu + \Psi_{3} + \Phi_{21}. \tag{3.18}
\]
\[
\Delta \nu - \bar{\delta}\gamma = 2\alpha \nu + (\bar{\gamma} - 3\gamma - \mu - \bar{\mu})\lambda - \Psi_{4}, \tag{3.19}
\]

\[
\delta \rho - \bar{\delta}\sigma = (\bar{\beta} + \alpha)\rho + (\bar{\beta} - 3\alpha)\sigma - \Psi_{1} + \Phi_{01}, \tag{3.20}
\]

\[
\delta \alpha - \bar{\delta}\beta = \mu\rho - \lambda\sigma - 2\alpha\beta + \alpha\bar{\alpha} + \beta\bar{\beta} - \Psi_{2} + \Lambda + \Phi_{11}, \tag{3.21}
\]

\[
\delta \lambda - \bar{\delta}\mu = (\alpha + \bar{\beta})\mu + (\bar{\alpha} - 3\beta)\lambda - \Psi_{3} + \Phi_{21}, \tag{3.22}
\]

\[
\delta \nu - \Delta \mu = \gamma\mu - 2\beta\tau + \bar{\gamma}\pi + \mu^{2} + \lambda\bar{\lambda} + \Phi_{22}, \tag{3.23}
\]

\[
\delta \gamma - \Delta \beta = \gamma\mu - \sigma\nu - (\mu - \gamma - \bar{\gamma})\beta + \lambda\alpha + \Phi_{12}, \tag{3.24}
\]

\[
\delta \tau - \Delta \sigma = 2\beta\tau + (\bar{\gamma} + \mu - 3\gamma)\sigma - \bar{\lambda}\rho + \Phi_{02}, \tag{3.25}
\]

\[
\Delta \rho - \bar{\delta}\tau = (\gamma + \bar{\gamma} - \bar{\mu})\rho - 2\alpha\tau - \lambda\sigma - \Psi_{2} - 2\Lambda, \tag{3.26}
\]

\[
\Delta \alpha - \delta \gamma = \rho\nu - \tau\lambda - \lambda\beta + (\bar{\gamma} - \gamma - \bar{\mu})\alpha - \Psi_{3}. \tag{3.27}
\]
\begin{align}
D &= L^{\mu} \nabla_{\mu} = \frac{\partial}{\partial r} \tag{3.28} \\[6pt]
\Delta &= N^{\mu} \nabla_{\mu} = V \frac{\partial}{\partial r} 
   + \frac{\partial}{\partial u} + X^{i}\frac{\partial}{\partial x^{i}} \tag{3.29} \\[6pt]
\delta &= M^{\mu} \nabla_{\mu} = W \frac{\partial}{\partial r} 
   + \xi^{i}\frac{\partial}{\partial x^{i}} \tag{3.30} \\[6pt]
\Phi_{00} &= -\tfrac{1}{2} R_{11} = \overline{\Phi}_{00} \tag{3.31} \\[6pt]
\Phi_{11} &= -\tfrac{1}{4}(R_{12} + R_{34}) = \overline{\Phi}_{11} \tag{3.32} \\[6pt]
\Phi_{01} &= -\tfrac{1}{2} R_{13} = \overline{\Phi}_{10} \tag{3.33} \\[6pt]
\Phi_{12} &= -\tfrac{1}{2} R_{23} = \overline{\Phi}_{21} \tag{3.34} \\[6pt]
\Phi_{02} &= -\tfrac{1}{2} R_{33} = \overline{\Phi}_{20} \tag{3.35} \\[6pt]
\Phi_{22} &= -\tfrac{1}{2} R_{22} = \overline{\Phi}_{22} \tag{3.36} \\[6pt]
\Lambda &= \tfrac{R}{24} \tag{3.37}
\end{align}

\begin{align}
\Psi_{0} &= -C_{1313} 
= - C_{\alpha \beta \gamma \delta} \, l^{\alpha} m^{\beta} l^{\gamma} m^{\delta} 
\tag{3.38} \\[10pt]
%
\Psi_{1} &= -C_{1213} 
= - C_{\alpha \beta \gamma \delta} \, l^{\alpha} n^{\beta} l^{\gamma} m^{\delta} 
\tag{3.39} \\[10pt]
%
\Psi_{2} &= -\tfrac{1}{2} \left( C_{1212} + C_{1234} \right) 
= -\tfrac{1}{2} C_{\alpha \beta \gamma \delta} 
\left( l^{\alpha} n^{\beta} l^{\gamma} n^{\delta} 
+ l^{\alpha} n^{\beta} m^{\gamma} \bar{m}^{\delta} \right) 
\tag{3.40} \\[10pt]
%
\Psi_{3} &= C_{1224} 
= C_{\alpha \beta \gamma \delta} \, l^{\alpha} n^{\beta} \bar{m}^{\gamma} n^{\delta} 
\tag{3.41} \\[10pt]
%
\Psi_{4} &= -C_{2424} 
= - C_{\alpha \beta \gamma \delta} \, n^{\alpha} \bar{m}^{\beta} n^{\gamma} \bar{m}^{\delta} 
\tag{3.42}
\end{align}





%%%%

Expressing the rotation coefficients in terms of the metric, we have:

\begin{align}
D \, \xi^{i} &= \rho \, \xi^{i} + \sigma \, \bar{\xi}^{i} 
\tag{3.43} \\[10pt]
%
D \, \omega &= \rho \, \omega + \sigma \, \bar{\omega} - (\alpha + \bar{\beta}) 
\tag{3.44} \\[10pt]
%
D \, \chi &= \tau \, \xi^{i} + \bar{\tau} \, \bar{\xi}^{i} 
\tag{3.45} \\[10pt]
%
D \, U &= \tau \, \omega + \bar{\tau} \, \bar{\omega} - (\gamma + \bar{\gamma}) 
\tag{3.46} \\[10pt]
%
\delta \chi^{i} - \Delta \xi^{i} &= (\mu + \bar{\gamma} - \gamma)\, \xi^{i} + \lambda \, \bar{\xi}^{i} 
\tag{3.47} \\[10pt]
%
\delta \, \bar{\xi}^{i} - \bar{\delta} \, \xi^{i} &= (\bar{\beta} - \alpha)\, \xi^{i} + (\bar{\alpha} - \beta)\, \bar{\xi}^{i} 
\tag{3.48} \\[10pt]
%
\delta \, \bar{\omega} - \bar{\delta} \, \omega &= (\bar{\beta} - \alpha)\, \omega + (\bar{\alpha} - \beta)\, \bar{\omega} + (\mu - \bar{\mu}) 
\tag{3.49} \\[10pt]
%
\delta U - \Delta \omega &= (\mu + \bar{\gamma} - \gamma)\, \omega + \bar{\lambda} \, \bar{\omega} - \nu 
\tag{3.50}
\end{align}



As in Chapter 2 we use the Bianchi identities as field equations for the Weyl tensor.  
In the Newman–Penrose formalism they may be written:  
(I am indebted to \textbf{R. G. McLenaghan} for these)

\begin{align}
\bar{\delta} \Psi_{0} - D \Psi_{1} + D \Phi_{01} - \delta \Phi_{00} 
&= (\pi - 4 \alpha) \Psi_{0} - 2 (2 \tau + \beta) \Phi_{00} \notag \\
& \quad + 2 \rho \Phi_{01} - 2 \delta \Phi_{10}
\tag{3.51} \\[12pt]
%
\Delta \Psi_{0} - \delta \Psi_{1} + D \Phi_{02} - \delta \Phi_{01} 
&= (4 \gamma - \mu) \Psi_{0} - 2 (2 \tau + \beta) \Psi_{1} \notag \\
& \quad + 3 \sigma \Psi_{2} - \lambda \Phi_{00} - 2 \beta \Phi_{01} + 2 \delta \Phi_{11} - \gamma \Phi_{02}
\tag{3.52} \\[12pt]
%
3 (\bar{\delta} \Psi_{1} - D \Psi_{2}) + 2 (D \Phi_{11} - \delta \Phi_{10}) 
&+ \bar{\delta} \Phi_{01} - \Delta \Phi_{00} \notag \\
&= 3 \lambda \Psi_{0} - 9 \rho \Psi_{2} + 6 \alpha \Psi_{1} \notag \\
& \quad + (\mu - 2 \gamma - 2 \bar{\gamma}) \Phi_{00} + (2 \alpha + 2 \bar{\pi}) \Phi_{01} \notag \\
& \quad + 2 (\tau - 2 \alpha) \Phi_{12} + 2 \rho \Phi_{11} + 2 \delta \Phi_{20} - 6 \Phi_{02}
\tag{3.53} \\[12pt]
%
3 (\Delta \Psi_{1} - \delta \Psi_{2}) + 2 (\Delta \Phi_{01} - \delta \Phi_{11}) 
&+ (\delta \Phi_{02} - \Delta \Phi_{01}) \notag \\
&= 3 \nu \Psi_{0} + 6 (\gamma - \mu) \Psi_{1} - 9 \tau \Psi_{2} + 6 \sigma \Psi_{3} \notag \\
& \quad + 2 (\mu - \bar{\mu} - \gamma) \Phi_{01} - 2 \lambda \Phi_{10} + 2 \nu \Phi_{11} \notag \\
& \quad + (2 \alpha + \bar{\pi} - 2 \beta) \Phi_{02} + 2 \delta \Phi_{12}
\tag{3.54} \\[12pt]
%
3 (\bar{\delta} \Psi_{2} - D \Psi_{3}) + 2 (D \Phi_{21} - \delta \Phi_{20}) 
&+ 2 (D \Phi_{11} - \delta \Phi_{21}) \notag \\
&= 6 \mu \Psi_{1} - 6 \rho \Psi_{3} - 2 \pi \Phi_{02} + 2 (\mu - \bar{\mu} - \gamma) \Phi_{10} \notag \\
& \quad + 4 \bar{\alpha} \Phi_{11} + (2 \beta - 2 \alpha - 2 \pi) \Phi_{20} - 2 \delta \Phi_{12}
\tag{3.55}
\end{align}

\begin{align}
3(\Delta \Psi_{2} - \delta \Psi_{3}) + (D \Phi_{22} - \delta \Phi_{21}) + 2 (\bar{\delta} \Phi_{12} - \Delta \Phi_{11}) 
&= 6 \nu \Psi_{2} - 9 \mu \Psi_{2} + 6 (\beta - \tau) \Psi_{3} - 3 \sigma \Psi_{4} \notag \\
& \quad - 2 \nu \Phi_{01} + 2 \bar{\nu} \Phi_{10} + 2 (\bar{\mu} - \mu) \Phi_{11} + 2 \lambda \Phi_{02} \notag \\
& \quad - \lambda \Phi_{20} + 2 (\bar{\pi} - 2 \beta) \Phi_{12} + 2 (\beta \tau - \gamma \mu) \Phi_{21} - 6 \Phi_{22}
\tag{3.56} \\[12pt]
%
\bar{\delta} \Psi_{3} - D \Psi_{4} + \bar{\delta} \Phi_{21} - \Delta \Phi_{20} 
&= 3 \lambda \Psi_{2} - 2 \nu \Psi_{3} - \rho \Psi_{4} \notag \\
& \quad - 2 \nu \Phi_{01} + 2 \lambda \Phi_{10} + (2 \gamma - 2 \bar{\gamma} - \mu) \Phi_{20} \notag \\
& \quad + 2 (\bar{\pi} - \alpha) \Phi_{21} - 6 \Phi_{22}
\tag{3.57} \\[12pt]
%
\Delta \Psi_{3} - \delta \Psi_{4} + \bar{\delta} \Phi_{22} - \Delta \Phi_{21} 
&= 3 \nu \Psi_{2} - 2 (\gamma + 2 \mu) \Psi_{3} + (4 \beta - \tau) \Psi_{4} \notag \\
& \quad - 2 \nu \Phi_{11} - \bar{\nu} \Phi_{00} + 2 \lambda \Phi_{02} \notag \\
& \quad + 2 (\gamma + \bar{\mu}) \Phi_{11} + (\bar{\pi} - 2 \beta - 2 \tau) \Phi_{22}
\tag{3.58} \\[12pt]
%
D \Phi_{12} - \delta \Phi_{11} - \bar{\delta} \Phi_{02} + \Delta \Phi_{01} + 3 \Lambda 
&= (2 \gamma - \mu - 2 \bar{\mu}) \Phi_{01} + \pi \bar{\nu} \Phi_{00} \notag \\
& \quad - \lambda \Phi_{02} - 2 \pi \Phi_{11} + (2 \beta - 2 \alpha - \bar{\pi}) \Phi_{02} + 3 \rho \Phi_{22} + 6 \Phi_{21}
\tag{3.59} \\[12pt]
%
D \Phi_{10} - \delta \Phi_{01} + \Delta \Phi_{00} + 3 D \Lambda 
&= (2 \gamma - \mu + 2 \bar{\gamma} - \bar{\mu}) \Phi_{00} - 2 (\alpha + \bar{\pi}) \Phi_{10} \notag \\
& \quad - 2 (\tau + \alpha) \Phi_{10} + 4 \rho \Phi_{11} + 6 \Phi_{02} - \tau \Phi_{20}
\tag{3.60} \\[12pt]
%
D \Phi_{22} - \delta \Phi_{21} - \bar{\delta} \Phi_{12} + \Delta \Phi_{11} + 3 \Delta \Lambda 
&= \nu \Phi_{01} + \bar{\nu} \Phi_{10} - 2 (\mu + \bar{\mu}) \Phi_{11} \notag \\
& \quad - \lambda \Phi_{02} - \bar{\lambda} \Phi_{20} + (2 \bar{\beta} - \tau) \Phi_{12} + (2 \beta - \tau) \Phi_{21} \notag \\
& \quad + 2 \rho \Phi_{22}
\tag{3.61}
\end{align}

\sectiontoc{4. The Undisturbed Metric}

The undisturbed metric may be written
\[
ds^{2} = \Omega^{2} \left( dt^{2} - d\rho^{2} - \sinh^{2} \rho \, ( d\theta^{2} + \sin^{2}\theta \, d\phi^{2} ) \right)
\]
where
\[
\Omega = A(\cosh t - 1).
\]

Put
\[
u = t - \rho,
\]
then
\[
ds^{2} = \Omega^{2} \left[ -du^{2} + 2 \, du \, dt - \sinh^{2}(t-u)\,(d\theta^{2} + \sin^{2}\theta \, d\phi^{2}) \right]
\tag{4.1}
\]

$u$ is a null coordinate.  

To calculate $r$, the affine parameter, we note that $t$ is an affine parameter for the metric within the square brackets.  
Therefore
\[
r = \int \Omega^{2} \, dt = B(u,\theta,\phi)
\tag{4.2}
\]
will be an affine parameter for (4.1).  

$B$ is constant along the null geodesic. Normally it would be taken so that $r=0$ when $t=u$.  
However, in our case it will be more convenient to make it zero and define $r$ as
\[
r = \int_{0}^{t} \Omega^{2} \, dt'
\tag{4.3}
\]

This means that surfaces of constant $r$ are surfaces of constant $t$.  
This may seem rather odd, but it should be pointed out that the choice of $B$ will not affect the asymptotic dependence of quantities.  
That is, if
\[
f = O(r^{-n})
\]
then
\[
f = O(r'^{-n}), \qquad r' = r + B.
\]


It proves easier to perform the calculations with this choice of $r$  
but all results could be transformed back to a more normal coordinate system.  

From (4.3)
\[
r = A^{2} \left[ \tfrac{1}{4} \sinh 2t - 2 \sinh t + \tfrac{3}{2} t \right]
\tag{4.4}
\]

The matter in the universe is assumed to be dust so its energy tensor may be written
\[
T_{ab} = \mu \, V_{a} V_{b}
\tag{4.5}
\]

For the undisturbed case, from Chapter 2
\[
\mu = \frac{6A}{\Omega^{3}}
\tag{4.6}
\]
\[
V_{a} = \Omega \, t_{;a}, 
\qquad 
V_{a} V^{a} = 1
\]

Now
\[
\Omega = \sqrt{2s} + A - \frac{3A^{2}}{2\sqrt{2}} \, \frac{\log s}{s} + O(s^{-1})
\tag{4.7}
\]
where
\[
s^{2} = r.
\]

Therefore if we try to expand $\mu$ as a series in powers of $s$  
the result will be very messy and will involve terms of the form
\[
\frac{\log^{n} s}{s^{n}} x^{(*)} 
\]

\bigskip
\noindent\rule{\linewidth}{0.4pt}  

\textit{* It should be pointed out that the expansions used will only be assumed to be valid asymptotically. They will not be assumed to converge at finite distances nor will the quantities concerned be assumed analytic. (see A. Erdelyi: \textit{Asymptotic Expansions} – Dover)}  \\
\bigskip
\noindent\rule{\linewidth}{0.4pt}  

This does not invalidate it as an asymptotic expansion but it makes it tedious to handle.  
For convenience therefore, we will perform the expansions in terms of $\Omega(r)$ which will be defined in general as the same function of $r$ as it is in the undisturbed case.  
That is
\[
\Omega = A(\cosh t - 1),
\qquad
r = A^{2} \left[ \tfrac{1}{4} \sinh 2t - 2 \sinh t + \tfrac{3}{2} t \right]
\tag{4.8}
\]
then
\[
\frac{d\Omega}{d r} = \frac{\sqrt{1 + \tfrac{3}{2}A}}{\Omega}
= \frac{1}{\Omega} \left[ 1 + \frac{A}{\Omega} - \frac{A^{2}}{2 \Omega^{2}} + \frac{A^{3}}{2 \Omega^{3}} - \frac{5A^{4}}{8 \Omega^{4}} + \frac{7A^{5}}{8 \Omega^{5}} - \cdots \right]
\tag{4.9}
\]

For the third and fourth coordinates it is more convenient to use stereographic coordinates than spherical polars.  

Since the matter is dust its energy–momentum tensor and hence the Ricci tensor have only four independent components.  
We will take these as $\Lambda, \, \Phi_{00}, \, \Phi_{01}$ (since $\Phi_{01}$ is complex it represents two components).  

In terms of these the other components of the Ricci tensor may be expressed as:
\[
\Phi_{11} = 3 \Lambda + \frac{\Phi_{01} \bar{\Phi}_{01}}{\Phi_{00}}
\]
\[
\Phi_{22} = \frac{3}{\Phi_{00}} \Lambda^{2} \left( 1 + \frac{\Phi_{01} \bar{\Phi}_{01}}{6 \Lambda \Phi_{00}} \right)^{2}
\]

\[
\Phi_{12} = \bar{\Phi}_{21} = \frac{6 \Lambda \Phi_{01}}{\Phi_{00}} 
\left( 1 + \frac{\Phi_{01} \bar{\Phi}_{01}}{6 \Lambda \Phi_{00}} \right)^{2}
\]

\[
\Phi_{02} = \bar{\Phi}_{20} = \frac{\Phi_{01}^{2}}{\Phi_{00}}
\tag{4.10}
\]

---

For the undisturbed universe with the coordinate system given:

\[
\Lambda = \frac{\mu}{24} = \frac{A}{4 \Omega^{3}}
\]

\[
\Phi_{00} = \frac{3A}{\Omega^{5}}, 
\qquad 
\Phi_{11} = \frac{3A}{4 \Omega^{3}}, 
\qquad 
\Phi_{22} = \frac{3A}{\Omega}
\]

\[
\Phi_{01} = \Phi_{02} = 0
\tag{4.11}
\]

---

Using these values and the fact that in the undisturbed universe all the $\Psi^{(n)}$ are zero,  
we may integrate equations (3.10–50) to find the values of the spin coefficients for the unperturbed universe:

\[
\rho = -\frac{2}{\Omega^{2}} - \frac{A}{\Omega^{3}} 
+ \left( \frac{A^{2}}{2} - \frac{A^{2}}{2} e^{2u} \right) \Omega^{-4}
+ A^{3} (u - e^{2u}) \Omega^{-5} 
+ A^{4} \left( \tfrac{5}{8} - \tfrac{1}{4}u - \tfrac{1}{8} e^{4u} \right) \Omega^{-6} + \cdots
\]

\[
\epsilon = \pi = \kappa = \nu = \lambda = \chi = 0
\]

\[
\xi^{2} = \frac{S e^{iu} - A S e^{iu}}{\Omega^{2}} 
= \frac{S e^{iu}}{\Omega^{3}} 
+ \tau A^{4} (5 S A u + S e^{3u}) \Omega^{-4} + \cdots
\]



\begin{align}
\alpha
&= -\bar{\beta}
 = \frac{i}{2}\frac{\bar{V}\bar{S}e^{i\alpha}}{\Omega^{2}}
 = \frac{i}{2}\frac{A\bar{V}\bar{S}e^{i\alpha}}{\Omega^{2}}
 + \cdots
\\[4pt]
\mu
&= \frac{A}{2\Omega}
 - \frac{\Omega^{2}}{4}\left(1+e^{2u}\right)\Omega^{-2}
 + \frac{\Omega^{3}}{4}\left(1+2e^{2u}\right)\Omega^{-3}
 + \cdots
\\[4pt]
\gamma
&= -\frac{1}{2}
 - \frac{A}{2\Omega}
 + \frac{A^{2}}{4\Omega^{2}}
 - \frac{A^{3}}{4\Omega^{3}}
 + \cdots
\\[4pt]
U
&= \frac{i}{2}\Omega^{2}
\tag{4.12}
\end{align}

\section{Boundary Conditions}

We wish to consider radiation in a universe that
asymptotically approaches the undisturbed universe given
above. $\phi_{00}$ and $\Lambda$ will then have the values
given above plus terms of smaller order. To determine this
order and the order of $\phi_{01}$ and $\gamma_0$, there are two
ways in which we may proceed. We may take the smallest orders
that will permit radiation, that is
\[
    \psi_4 = O(r^{-1}).
\]

Larger order terms than these in $\phi_{00}$, $\Lambda$ and
$\phi_{01}$ turn out to have their $u$ derivatives
dependent only on themselves and not on the $r^{-1}$ coefficient
of $\psi_4$, the radiation field. They are thus disturbances
not produced by the radiation field and will not be considered.

Alternatively we may proceed by a method of successive approximations. We take the undisturbed values of the spin
coefficients and use them to solve the Bianchi Identities as
field equations for the conformal tensor using the flat space
boundary condition that
\[
\psi_0 = O(r^{-5}).
\]
Then substituting these $\psi$'s in equations (3.10--27) calculate
the disturbances induced in the spin coefficients and substituting
these back in the Bianchi Identities, calculate the disturbances in
the $\psi$'s. Further iteration does not affect the orders
of the disturbances.

Both these methods indicate that the boundary conditions
should be:
\begin{align}
\Lambda
    &= \frac{A}{4\Omega^3} + O(\Omega^{-7})
    \tag{5.1}\\[6pt]
\phi_{00}
    &= \frac{3A}{\Omega^5} + O(\Omega^{-9})
    \tag{5.2}\\[6pt]
\phi_{01}
    &= O(\Omega^{-7})
    \qquad \text{(see next section)}
    \tag{5.3}\\[6pt]
\gamma_0
    &= O(\Omega^{-7})
    \tag{5.4}
\end{align}

We also assume ``uniform smoothness'', that is:
\[
\frac{\partial}{\partial x^i}
\frac{\partial}{\partial x^j}
\cdots
\frac{\partial}{\partial x^k}
\Lambda
=
O(\Omega^{-7}).
\]

\[
\frac{\partial \Lambda}{\partial \Omega}
=
-\frac{3A}{4\Omega^4}
+ O(\Omega^{-8}).
\]

\hfill etc.\ldots/


It will be shown that if these boundary conditions
hold on one hypersurface ($u = \mathrm{const.}$) they will hold on
succeeding hypersurfaces and that these conditions are the
most severe to permit radiation.

\section{Integration}

As Newman and Penrose, we begin by integrating the
equations (3.10 \& 11)
\begin{align}
D\rho &= \rho^2 + \sigma\bar{\sigma} + \phi_{00},\\
D\sigma &= 2\rho\sigma + \psi_0.
\end{align}

where
\[
\phi_{00}
=
\frac{3A}{4\sqrt{2}\,r^{3/2}}
+ O(r^{-3})
\]

\[
\psi_0
=
\frac{\psi_0^{\,0}}
     {8\sqrt{2}\,r^{7/2}}
+ O(r^{-4}).
\]

Let
\[
P =
\begin{bmatrix}
\rho & \sigma\\
\bar{\sigma} & \rho
\end{bmatrix},
\qquad
\Phi =
\begin{bmatrix}
\phi_{00} & \psi_0\\
\bar{\psi}_0 & \phi_{00}
\end{bmatrix}.
\]


then
\begin{equation}
D P = P^2 + \Phi
\tag{6.1}
\end{equation}

let
\begin{equation}
P = -(D Y)Y^{-1}
\tag{6.2}
\end{equation}

then
\begin{equation}
D^2 Y = -\Phi Y
\tag{6.3}
\end{equation}

since
\[
\int r\Phi\,dr < \infty
\]

\[
DY = F + O(1),
\qquad
\text{where $F$ is constant}
\tag{6.4}
\]

\[
Y = rF + O(r).
\]

However
\[
\Phi = O(r^{-5/2}).
\]

therefore
\begin{equation}
D^2Y
=
-r\Phi F + O(r^{-3/2})
=
O(r^{-3/2})
\tag{6.5}
\end{equation}

therefore
\[
DY = F + O(r^{-1/2})
\]

\[
Y = rF + O(r^{1/2}) + E,
\qquad
\text{$E$ is constant}
\]

\begin{equation}
P
=
-r^{-1}I + O(r^{-3/2})
\tag{6.6}
\end{equation}

if $F$ is non-singular (The case $F$ singular corresponds to
asymptotically plane or cylindrical surfaces and will not be
considered here).

Thus
\[
\rho
=
-r^{-1} + O(r^{-3/2})
=
-2\Omega^{-2} + O(\Omega^{-3}),
\]

\begin{equation}
\sigma
=
O(r^{-3/2})
=
O(\Omega^{-3})
\tag{6.7}
\end{equation}

Let
\begin{equation}
\rho
=
-2\Omega^{-2} + g\Omega^{-3}
\tag{6.8}
\end{equation}

\[
\sigma = h\Omega^{-3},
\]

where
\[
g,h,\ldots = O(1).
\]

Then using:
\[
D
=
\frac{\partial}{\partial r}
=
-\Omega^{-1}
\left(
1 + A\Omega^{-1}
-\frac{A^2\Omega^2}{2}
+\frac{A^3\Omega^3}{2}
-\cdots
\right)
\frac{\partial}{\partial\Omega}.
\]
\[
\frac{\partial g}{\partial \Omega}
\left(\Omega + O(1)\right)
=
-A-g+O(\Omega^{-1})
\]

Integrating,
\[
g
=
\frac{1}{\Omega+O(1)}
\left[
a+
\int
\frac{
\left(-A+O(\Omega^{-1})\right)
\left(\Omega+O(1)\right)
}{
\Omega+O(1)
}
\,d\Omega
\right].
\]

therefore
\begin{equation}
g
=
-A
+
O\left(\frac{\log\Omega}{\Omega}\right)
\tag{6.9}
\end{equation}

for $h$,
\[
\frac{\partial h}{\partial\Omega}
\left(\Omega+O(1)\right)
=
-h+O(\Omega^{-1})
\]

therefore
\begin{equation}
h
=
O(\Omega^{-1})
\tag{6.10}
\end{equation}

Repeat the process with
\begin{align}
\rho &= -2\Omega^{-2} - A\Omega^{-3} + g\Omega^{-4},
\tag{6.12}\\
\sigma &= h\Omega^{-4}.
\end{align}

where
\[
g = O(\log\Omega),
\]
\[
h = O(1).
\]

then
\[
\frac{\partial h}{\partial\Omega}
\left(\Omega + O(1)\right)
=
O(\Omega^{-1})
\]

\begin{equation}
h
=
\sigma^{0}(u,x^{c})
+
O(\Omega^{-1})
\tag{6.13}
\end{equation}

\[
\frac{\partial g}{\partial\Omega}
\left(-\Omega + O(1)\right)
=
O(\Omega^{-1}\log\Omega)
\]

\begin{equation}
g
=
\rho^{0}(u,x^{c})
+
O(\Omega^{-1}\log\Omega)
\tag{6.14}
\end{equation}

Unlike Newman and Unti, we cannot make $\rho^{0}$ zero by
the transformation
\[
r' = r - \rho^{0},
\]
since this would alter the boundary condition
\[
\Lambda
=
\frac{A}{4\Omega^{3}}
+
\frac{\Lambda^{0}}{\Omega^{7}}.
\]

Continuing the above process we derive:
\begin{equation}
\begin{aligned}
\rho
={}&
-2\Omega^{-2}
-A\Omega^{-3}
+\rho^{0}\Omega^{-4}
+\left(
\frac{1}{2}A^{2}
-2A\rho^{0}
\right)\Omega^{-5}
\\
&+
\left(
-\frac{5}{4}A^{4}
+4A^{2}\rho^{0}
-\frac{(\rho^{0})^{2}}{2}
-\frac{\sigma^{0}\bar{\sigma}^{0}}{2}
\right)\Omega^{-6}
+O(\Omega^{-7})
\end{aligned}
\tag{6.15}
\end{equation}

\begin{equation}
\sigma
=
\sigma^{0}\Omega^{-4}
-
\left(
2A\sigma^{0}
+
\psi_{0}^{0}
\right)\Omega^{-5}
+\cdots
\tag{6.16}
\end{equation}



To determine the asymptotic behaviour of
$\psi_1$, $\alpha$, $\beta$, $\xi^c$ and $\omega$
we use the lemma proved by Newman and Penrose:

The $n\times n$ matrix $B$ and the column vector $b$ are
given functions of $x$ such that:
\begin{equation}
B = O(x^{-2}),
\qquad
b = O(x^{-2})
\tag{6.17}
\end{equation}

The $n\times n$ matrix $A$ is independent of $x$ and has no
eigenvalue with positive real part. Any eigenvalue with
vanishing real part is regular. Then all solutions of:
\begin{equation}
\frac{\partial}{\partial x}y
=
\left(Ax^{-1}+B\right)y+b
\tag{6.18}
\end{equation}
are bounded as $x\rightarrow\infty$. $y$ is a column
vector.

For reasons to be explained below, we will assume for
the moment that
\[
\phi_{01}=O(\Omega^{-5}),
\]

\[
\frac{\partial}{\partial\Omega}\phi_{01}
=
O(\Omega^{-6}),
\tag{6.19}
\]

\[
\frac{\partial}{\partial x^c}
\frac{\partial}{\partial x^d}
\cdots
\phi_{01}
=
O(\Omega^{-5}).
\]

We take as $y$ the column vector
\begin{equation}
y=
\left[
\Omega^5\psi_1,\,
\Omega^2\alpha,\,
\Omega^2\bar{\alpha},\,
\Omega^2\beta,\,
\Omega^2\bar{\beta},\,
\Omega^2\xi^3,\,
\Omega^2\bar{\xi}^3,\,
\Omega^2\xi^4,\,
\Omega^2\bar{\xi}^4,\,
\omega,\,
\bar{\omega}
\right]^T .
\tag{6.20}
\end{equation}

By equations 3.51, 3.13, 3.14, 3.43, 3.44,

\begin{equation}
A=
\begin{pmatrix}
-3 & 0 & 6A & 6A & 0 & 0 & 0 & 0 & 0 & 1.5A & 0 \\
 0 & 0 & 0  & 0  & 0 & 0 & 0 & 0 & 0 & 0    & 0 \\
 0 & 0 & 0  & 0  & 0 & 0 & 0 & 0 & 0 & 0    & 0 \\
 0 & 0 & 0  & 0  & 0 & 0 & 0 & 0 & 0 & 0    & 0 \\
 0 & 0 & 0  & 0  & 0 & 0 & 0 & 0 & 0 & 0    & 0 \\
 0 & 0 & 0  & 0  & 0 & 0 & 0 & 0 & 0 & 0    & 0 \\
 0 & 0 & 0  & 0  & 0 & 0 & 0 & 0 & 0 & 0    & 0 \\
 0 & 0 & 0  & 0  & 0 & 0 & 0 & 0 & 0 & 0    & 0 \\
 0 & 0 & 0  & 0  & 0 & 0 & 0 & 0 & 0 & 0    & 0 \\
 0 & 0 & -1 & -1 & 0 & 0 & 0 & 0 & 0 & -2   & 0 \\
 0 & -1 & 0 & 0 & -1 & 0 & 0 & 0 & 0 & 0    & -2
\end{pmatrix}.
\tag{6.21}
\end{equation}

$B$ and $b$ are $O(\Omega^{-1})$ expressions involving
\[
\rho,\quad
\sigma,\quad
\psi_0,\quad
\frac{\partial\psi_0}{\partial\Omega},\quad
\frac{\partial\psi_0}{\partial x^c},\quad
\phi_{01},
\quad\text{and}\quad
\frac{\partial\phi_{01}}{\partial\Omega}.
\]

Thus
\begin{equation}
\begin{aligned}
\psi_1 &= O(\Omega^{-5}),\\
\alpha,\beta,\xi^3,\xi^4 &= O(\Omega^{-2}),\\
\omega &= O(1).
\end{aligned}
\tag{6.22}
\end{equation}

since
\[
\tau = \bar{\alpha}+\beta,
\]
\[
\tau = O(\Omega^{-2}).
\]


Using this we integrate equation (3.12) by the same
method as above. We obtain
\begin{equation}
\tau
=
\tau^{0}(u,x^{i})\Omega^{-2}
+
O(\Omega^{-3})
\tag{6.23}
\end{equation}

We may make a null rotation of the tetrad on each null
geodesic
\begin{equation}
\begin{aligned}
l'_{\mu} &= l_{\mu},\\
n'_{\mu} &= n_{\mu}
          + a\bar{m}_{\mu}
          + \bar{a}m_{\mu}
          + a\bar{a}\,l_{\mu},\\
m'_{\mu} &= m_{\mu}+a\,l_{\mu}.
\end{aligned}
\tag{6.24}
\end{equation}

$a$ is constant along the geodesic since the tetrad is
parallelly transported.

By taking
\[
a=-\frac{1}{2}\tau^{0}
\]
we may make
\[
\tau'^{0}=0.
\]

Under a null rotation
\[
\phi'_{01}
=
\phi_{01}
+
a\phi_{00}.
\]

Thus until we have specified the null rotation we cannot
impose a boundary condition on $\phi_{01}$ more severe than
\[
\phi_{01}=O(\Omega^{-5}).
\]

We will specify the null rotation by
\[
\tau^{0}=0
\]
and in that tetrad system will impose the boundary condition
that
\[
\phi_{01}=O(\Omega^{-7})
\]
and is uniformly smooth.

Then by using this condition on $\phi_{01}$ and
\[
\tau=O(\Omega^{-3}),
\]
by equation (3.44)
\[
\omega=O(\Omega^{-1}).
\]


using this in equation (3.51)
\[
\psi_1 = O(\Omega^{-6})
\]

then by equation (3.12)
\[
\tau = O(\Omega^{-4})
\]

putting this back in equation (3.44)
\[
\omega = O(\Omega^{-2}\log\Omega)
\]

by equation (3.51)
\[
\psi_1 = O(\Omega^{-7}\log\Omega)
\]

by equation (3.12)
\[
\tau = O(\Omega^{-5}\log\Omega)
\]

by equation (3.44)
\[
\omega
=
\omega^{0}\Omega^{-2}
+
O(\Omega^{-3}\log\Omega)
\]

by equation (3.51)
\begin{equation}
\psi_1 = O(\Omega^{-7})
\tag{6.25}
\end{equation}

by equation (3.12)
\begin{equation}
\tau = O(\Omega^{-5})
\tag{6.26}
\end{equation}

by equation (3.44)
\begin{equation}
\omega
=
\omega^{0}\Omega^{-2}
+
O(\Omega^{-3})
\tag{6.27}
\end{equation}

By differentiating the equations used with respect to $x^i$,
one may show that $\psi_1$, $\alpha$, $\beta$, $a$, $\xi^i$, $\omega$
are uniformly smooth.

Adding equations 3.53 and 3.60:
\begin{equation}
\begin{aligned}
\delta\psi_1
- D\psi_2
+ D\phi_{11}
- \delta\phi_{01}
+ D\Lambda
={}&
\lambda\psi_0
-3\rho\psi_2
+2\alpha\psi_1
-\mu\phi_{00}
-2\bar{\alpha}\phi_{10}
\\
&+2\rho\phi_{11}
+\sigma\phi_{20}.
\end{aligned}
\tag{6.28}
\end{equation}


we may use the lemma again with $y$

\[
y =
\begin{bmatrix}
\Omega^4 \psi_2 \\
\Omega^2 \lambda \\
\Omega \mu
\end{bmatrix}.
\]

By equations 3.16, 3.17, 6.28,

\[
A =
\begin{bmatrix}
-2 & 0 & 3A \\
0  & 0 & 0  \\
0  & 0 & -1
\end{bmatrix}.
\]

$B$ and $b$ are
\[
O(\Omega^{-2}).
\]

Therefore
\begin{equation}
\begin{aligned}
\psi_2 &= O(\Omega^{-4}),\\
\lambda &= O(\Omega^{-2}),\\
\mu &= O(\Omega^{-1}),
\end{aligned}
\tag{6.29}
\end{equation}
and are uniformly smooth.

From (6.29) and (3.17), we may show
\[
\mu
=
\frac{A}{2}\Omega^{-1}
+
O(\Omega^{-2}\log\Omega).
\]

using this in (6.28)
\[
\psi_2
=
O(\Omega^{-5}\log\Omega).
\]

then by (3.17)
\begin{equation}
\mu
=
\frac{A}{2}\Omega^{-1}
+
\mu^0(u,x^i)\Omega^{-2}
+
O(\Omega^{-3}\log\Omega).
\tag{6.30}
\end{equation}

and
\[
\psi_2 = O(\Omega^{-5}).
\]


Integrating the radial equations 3.13, 3.14, 3.15, 3.16,
3.43, 3.45, 3.46,

\begin{equation}
\alpha
=
\alpha^{0}\Omega^{-2}
-
A\alpha^{0}\Omega^{-3}
+
\frac{1}{2}
\left(
3A^{2}\alpha^{0}
-
\rho^{0}\alpha^{0}
+
\bar{\sigma}^{0}\beta^{0}
\right)\Omega^{-4}
+
O(\Omega^{-5})
\tag{6.31}
\end{equation}

\begin{equation}
\beta
=
\beta^{0}\Omega^{-2}
-
A\beta^{0}\Omega^{-3}
+
\frac{1}{4}
\left(
3A^{2}\beta^{0}
-
\rho^{0}\beta^{0}
+
\bar{\beta}^{0}\sigma^{0}
\right)\Omega^{-4}
+
O(\Omega^{-5})
\tag{6.32}
\end{equation}

\begin{equation}
\gamma
=
\gamma^{0}
-
\frac{A}{2}\Omega^{-1}
+
\frac{A^{2}}{4}\Omega^{-2}
+
O(\Omega^{-3})
\tag{6.33}
\end{equation}

\begin{equation}
\lambda
=
\lambda^{0}\Omega^{-2}
-
A\left(
\lambda^{0}
+
\frac{\bar{\sigma}^{0}}{2}
\right)\Omega^{-3}
+
O(\Omega^{-4})
\tag{6.34}
\end{equation}

\begin{equation}
\xi^{c}
=
\xi^{0c}\Omega^{-2}
+
A\Omega^{3}\xi^{0c}
-
\frac{1}{2}
\left(
3A^{2}\xi^{0c}
-
\rho^{0}\xi^{0c}
-
\bar{\xi}^{0c}\sigma^{0}
\right)\Omega^{-4}
+
O(\Omega^{-5})
\tag{6.35}
\end{equation}

\begin{equation}
X^{c}
=
X^{0c}
+
O(\Omega^{-5})
\tag{6.36}
\end{equation}

\begin{equation}
\begin{aligned}
U
={}&
-\frac{1}{2}
\left(
\gamma^{0}
+
\bar{\gamma}^{0}
\right)\Omega^{2}
-
A\left(
\gamma^{0}
+
\bar{\gamma}^{0}
\right)\Omega
-
\frac{3}{2}A^{2}
\\
&\quad\times
\left(
1+
\gamma^{0}
+
\bar{\gamma}^{0}
\right)\log\Omega
+
U^{0}
+
O(\Omega^{-1})
\end{aligned}
\tag{6.37}
\end{equation}

Adding equation (3.55) $+\,2\times\overline{(3.59)}$,

\begin{equation}
\begin{aligned}
\delta\psi_{2}
-
D\psi_{3}
+
D\phi_{21}
-
\delta\phi_{20}
+
2\delta\Lambda
={}&
2\lambda\psi_{1}
-
2\rho\psi_{3}
\\
&-
\frac{1}{3}
\left(
\bar{\mu}+2\mu
\right)\phi_{10}
+
2(\beta-\bar{\alpha})\phi_{20}
+
2\rho\phi_{21}.
\end{aligned}
\tag{6.38}
\end{equation}


Therefore
\begin{equation}
\psi_3
=
\psi_3^{0}\Omega^{-4}
+
O(\Omega^{-5})
\tag{6.39}
\end{equation}

By equation (3.18)
\begin{equation}
\nu
=
\nu^{0}
-
\frac{1}{2}\psi_3^{0}\Omega^{-2}
+
O(\Omega^{-3})
\tag{6.40}
\end{equation}

By the orthonormality relations (2.2),
\[
g^{ui}
=
X^i
-
\left(
\xi^i\bar{\omega}
+
\bar{\xi}^i\omega
\right)
=
X^{i0}
+
O(\Omega^{-4}),
\]

\[
g^{ij}
=
-\left(
\xi^i\bar{\xi}^j
+
\bar{\xi}^i\xi^j
\right),
\qquad i,j=3,4,
\]

\begin{equation}
g^{ij}
=
-\left(
\xi^{i0}\bar{\xi}^{j0}
+
\bar{\xi}^{i0}\xi^{j0}
\right)
\left(
\Omega^{-4}
-
2A\Omega^{-5}
\right)
+
O(\Omega^{-6})
\tag{6.41}
\end{equation}

By making the coordinate transformation
\[
u' = u,
\]
\[
r' = r,
\]
\[
x'^{\,i}
=
x^i
+
C^i(u,x^c),
\]

where
\begin{equation}
\begin{aligned}
C^3{}_{,1}
&=
-X^{30}\left(1+C^3{}_{,3}\right)
-X^{40}C^3{}_{,4},
\\
C^4{}_{,1}
&=
-X^{40}\left(1+C^4{}_{,4}\right)
-X^{30}C^4{}_{,3},
\end{aligned}
\tag{6.42}
\end{equation}

\[
X^{i0}=0.
\]

We still have the coordinate freedom
\[
x'^i = D^i(x^j).
\]

We may use this to reduce the leading term of
$g^{ij}$ $(i,j=3,4)$ to a conformally flat metric
(c.f. Newman and Unti), that is:
\begin{equation}
g^{ij}
=
-2P\bar{P}\,\delta^{ij}
\left(
\Omega^{-4}
-
2A\Omega^{-5}
\right)
+
O(\Omega^{-6}).
\tag{6.43}
\end{equation}

where
\[
\xi^{30} = -i\xi^{40} = P(u,x^i).
\]

\section{Non-radial Equations}

By comparing coefficients of the various powers of $\Omega$
in the non-radial equations of \S 3, relations between
the integration constants of the radial equations may be
obtained:

In equation 3.23 the term in $\Omega^{-1}$ is
\[
-\frac{3}{4}A\left(\gamma^0+\bar{\gamma}^0\right)
-\frac{3}{4}A
\]

therefore
\begin{equation}
\gamma^0+\bar{\gamma}^0=-1.
\tag{7.1}
\end{equation}

therefore by (6.37)
\begin{equation}
U
=
\frac{1}{2}\Omega^2
+
U^0
+
O(\Omega^{-1}).
\tag{7.2}
\end{equation}

In equation (3.50), the constant term is
\[
\nu^0
\]

therefore
\begin{equation}
\nu^0=0.
\tag{7.3}
\end{equation}

By the $\Omega^{-2}$ term in (3.47)
\[
P_{,1}-P_{,1}
=
\left(\bar{\gamma}^0-\gamma^0\right)P.
\]

By the $\Omega^{-3}$ term
\[
P_{,1}-P_{,1}=0.
\]

therefore
\begin{equation}
\gamma^0=\bar{\gamma}^0=-\frac{1}{2}.
\tag{7.4}
\end{equation}

\[
P = S(x^i)e^{u}
\]

if
\[
S = |S|e^{i\phi}.
\]

By making a spatial rotation of the tetrad
\[
m'^{\mu}=e^{-i\phi}m^{\mu}
\]
we make $S$ real. We take
\[
S=\frac{1}{\sqrt{2}}
\left(
1+\frac{1}{4}\left[(x^3)^2+(x^4)^2\right]
\right),
\]
the stereographic project factor for a 2-sphere.

By the $\Omega^{-4}$ term in equation (3.48)
\begin{equation}
\begin{aligned}
\alpha^0
&=
\frac{1}{2}\,\bar{\nabla}\bar{S}\,e^{u},
\\
\beta^0
&=
-\frac{1}{2}\,\nabla S\,e^{u}.
\end{aligned}
\tag{7.5}
\end{equation}

where
\[
\nabla
=
\frac{\partial}{\partial x^3}
+
i\frac{\partial}{\partial x^4}.
\]

By the $\Omega^{-4}$ term in (3.21)
\[
\frac{1}{2}e^{2u}
\left(
S\nabla\bar{\nabla}S
+
S\bar{\nabla}\nabla S
\right)
=
-2\mu^0
-
\frac{A^2}{2}
(\nabla P)(\bar{\nabla}\bar{P}).
\]

\begin{equation}
\begin{aligned}
\mu^0
&=
-\frac{A}{4}
-\frac{S^2}{2}\,
\nabla\bar{\nabla}\log S\,e^{2u}
\\
&=
-\frac{A^2}{4}
\left[
1+e^{2u}
\right].
\end{aligned}
\tag{7.6}
\end{equation}

By (3.59), (3.60), (3.61)
\[
\frac{\partial}{\partial u}\phi_{0i}
=
O(\Omega^{-7}).
\]
\[
\frac{\partial}{\partial u}\phi_{00}
=
O(\Omega^{-7})
\]

\begin{equation}
\frac{\partial\Lambda}{\partial u}
=
H\Omega^{-5}
+
O(\Omega^{-6}\log\Omega)
\tag{7.7}
\end{equation}

where
\[
H
=
-\frac{A^2}{8}
+
\frac{3U^0}{4}
+
\frac{\rho^0}{4}
+
\frac{A^2e^{2u}}{8}.
\]

By making a coordinate transformation*
\begin{equation}
u'
=
u
-
\int_{u_0}^{u} H\,du''
\tag{7.8}
\end{equation}

\[
H'=0.
\]

therefore
\[
U^0
=
\frac{A^2}{6}
-
\frac{A^2e^{2u}}{6}
-
\frac{\rho^0}{3}.
\]

By the $\Omega^{-4}$ term in (3.26)
\begin{equation}
\rho^0{}_{,1}
+
4U^0
=
-\frac{A^2}{2}
\left(1+e^{2u}\right)
\tag{7.9}
\end{equation}

therefore
\[
\rho^0{}_{,1}
-
\frac{7}{3}\rho^0
+
\frac{7}{6}A^2
-
\frac{1}{6}A^2e^{2u}
=
0.
\]

\begin{equation}
\rho^0
=
\frac{A^2}{2}
\left(1-e^{2u}\right)
+
C(x^i)e^{\frac{7}{3}u}.
\tag{7.10}
\end{equation}

\bigskip
\hrule
\medskip

\noindent
*This transformation does not upset the boundary conditions
on the hypersurface
\[
u=u^0.
\]

\medskip
\hrule
By the $\Omega^{-4}$ term in (3.25)
\begin{equation}
\lambda^0
=
\frac{1}{2}
\left(
\bar{\sigma}^0{}_{,1}
-
\bar{\sigma}^0
\right)
\tag{7.11}
\end{equation}

By the $\Omega^{-4}$ term in (3.22)
\begin{equation}
\psi_3^0
=
\left(
\bar{\sigma}^0{}_{,1}
-
\bar{\sigma}^0
\right)e^u \nabla S
-
\frac{1}{2}e^u S
\nabla
\left(
\bar{\sigma}^0{}_{,1}
-
\bar{\sigma}^0
\right)
\tag{7.12}
\end{equation}

By the $\Omega^{-2}$ term in (3.50)
\[
2\omega^0{}_{,1}
-
\omega^0{}_{,1}
=
\frac{1}{2}\psi_3^0
\]

\begin{equation}
\omega^0
=
e^u
\left(
S\bar{\nabla}\bar{\sigma}^0
-
2\sigma^0\bar{\nabla}S
\right)
+
K(x^i)e^{2u}
\tag{7.13}
\end{equation}

By the $\Omega^{-6}$ term in (3.20)
\[
e^u S\nabla\rho^0
+
4\omega^0
-
e^u S\bar{\nabla}\bar{\sigma}^0
=
2\sigma^0\bar{\nabla}S\,e^u
\]

therefore
\begin{equation}
C=K=0
\tag{7.14}
\end{equation}

therefore
\begin{equation}
U^0=0
\tag{7.15}
\end{equation}

\begin{equation}
\rho^0
=
\frac{A^2}{2}
\left(
1-e^{2u}
\right)
\tag{7.16}
\end{equation}

Using (7.6), (7.16), in (6.28)
\[
\psi_2
=
O(\Omega^{-6}\log\Omega)
\]

Then by (3.17)
\begin{equation}
\begin{aligned}
\mu
={}&
\frac{A}{2}\Omega^{-1}
-
\frac{A^2}{4}
\left(
1+e^{2u}
\right)\Omega^{-2}
\\
&-
\frac{A^3}{2}
\left(
\frac{1}{2}+e^{2u}
\right)\Omega^{-3}
+
O(\Omega^{-4}\log\Omega)
\end{aligned}
\tag{7.17}
\end{equation}

Using this in (6.28)
\begin{equation}
\psi_2
=
O(\Omega^{-6})
\tag{7.18}
\end{equation}



By (3.51), (3.59), (3.60), (3.61)
\[
\frac{\partial}{\partial u}\phi_{01}
=
O(\Omega^{-7})
\]

\[
\frac{\partial}{\partial u}\phi_{00}
=
O(\Omega^{-9})
\]

\[
\frac{\partial}{\partial u}\Lambda
=
O(\Omega^{-7})
\]

\begin{equation}
\frac{\partial}{\partial u}\psi_0
=
O(\Omega^{-7})
\tag{7.19}
\end{equation}

Therefore if the boundary conditions (5.1--4) hold on
one null hypersurface, they will hold on succeeding hypersurfaces.

By (3.57)
\begin{equation}
\psi_4
=
O(\Omega^{-2})
\tag{7.20}
\end{equation}

The ``peeling off'' behaviour is therefore:
\begin{equation}
\begin{aligned}
\psi_4 &= O(r^{-1}),\\
\psi_3 &= O(r^{-2}),\\
\psi_2 &= O(r^{-3}),\\
\psi_1 &= O(r^{-7/2}),\\
\psi_0 &= O(r^{-7/2}).
\end{aligned}
\tag{7.21}
\end{equation}

As mentioned before, this asymptotic behaviour is independent
of the zero of $r$ and will hold for any affine parameter $r$.

To perform the remaining integrations we will assume
for definiteness:
\begin{equation}
\begin{aligned}
\phi_{01}
&=
\phi_{01}^{0}\Omega^{-7}
+
\phi_{01}^{1}\Omega^{-8}
+
O(\Omega^{-9}),
\\
\psi_0
&=
\psi_0^{0}\Omega^{-7}
+
\psi_0^{1}\Omega^{-8}
+
O(\Omega^{-9}),
\\
\Lambda
&=
\frac{A}{4}\Omega^{-3}
+
\Lambda^{0}\Omega^{-7}
+
O(\Omega^{-8}),
\\
\phi_{00}
&=
3A\Omega^{-5}
+
\phi_{00}^{0}\Omega^{-9}
+
O(\Omega^{-10}).
\end{aligned}
\tag{7.22}
\end{equation}

Then:
\begin{equation}
\begin{aligned}
\rho
={}&
-2\Omega^{-2}
-A\Omega^{-3}
+\frac{A^2}{2}\left(1-e^{2u}\right)\Omega^{-4}
\\
&+
\frac{A^3}{2}
\left(1-2e^{2u}\right)\Omega^{-5}
\\
&+
\left[
A^4
\left(
\frac{5}{8}
-\frac{7}{8}e^{2u}
-\frac{1}{8}e^{4u}
\right)
-\frac{\sigma^0\bar{\sigma}^0}{2}
\right]\Omega^{-6}
\\
&+
\frac{1}{3}
\left[
A^5
\left(
-\frac{37}{8}
+8e^{2u}
+8e^{4u}
+2e^{6u}
\right)
-\phi_{00}^{0}
\right.
\\
&\qquad\left.
+\bar{\sigma}^{0}
\left(
3A\sigma^{0}
+\psi_0^{0}
\right)
+
\sigma^{0}
\left(
3A\bar{\sigma}^{0}
+\bar{\psi}_0^{0}
\right)
\right]\Omega^{-7}
\\
&+
O(\Omega^{-8}).
\end{aligned}
\tag{7.23}
\end{equation}

\begin{equation}
\sigma
=
\sigma^{0}\Omega^{-4}
-
\left(
2A\sigma^{0}
+
\psi_0^{0}
\right)\Omega^{-5}
+
O(\Omega^{-6}).
\tag{7.24}
\end{equation}

\begin{equation}
\tau
=
-\frac{1}{3}
\left(
\phi_{01}^{0}
+
\psi_1^{0}
\right)\Omega^{-5}
+
O(\Omega^{-6}\log\Omega).
\tag{7.26}
\end{equation}

\begin{equation}
\begin{aligned}
\omega
={}&
\frac{e^u}{4}
\left(
S\bar{\nabla}\sigma^0
-
2\sigma^0\bar{\nabla}S
\right)\Omega^{-2}
\\
&-
\left[
\frac{Ae^u}{4}
\left(
S\bar{\nabla}\sigma^0
-
2\sigma^0\bar{\nabla}S
\right)
+
\frac{1}{3}
\left(
\phi_{01}^0+\psi_1^0
\right)
\right]\Omega^{-3}
\\
&+
O(\Omega^{-4}\log\Omega).
\end{aligned}
\tag{7.26}
\end{equation}

\begin{equation}
\begin{aligned}
\gamma
={}&
-\frac{1}{2}
-\frac{A}{2}\Omega^{-1}
+\frac{A^2}{4}\Omega^{-2}
-\frac{A^3}{4}\Omega^{-3}
\\
&+
\frac{1}{4}
\left(
\frac{5}{4}A^4
-\frac{\psi_0^0}{2}
\right)\Omega^{-4}
+
O(\Omega^{-5}).
\end{aligned}
\tag{7.27}
\end{equation}

\begin{equation}
\begin{aligned}
\mu
={}&
\frac{A}{2}\Omega^{-1}
-
\frac{A^2}{4}
\left(
1+e^{2u}
\right)\Omega^{-2}
+
\frac{A^3}{4}
\left(
1+2e^{2u}
\right)\Omega^{-3}
\\
&-
\frac{1}{2}
\left[
A^4
\left(
\frac{5}{8}
+\frac{7}{8}e^{2u}
+\frac{1}{8}e^{4u}
\right)
+
\frac{S^2}{2}
\left(
\bar{\sigma}^0{}_{,1}
-
\sigma^0{}_{,1}
\right)
+
\frac{\psi_0^0}{2}
\right]\Omega^{-4}
\\
&+
O(\Omega^{-5}).
\end{aligned}
\tag{7.28}
\end{equation}

\begin{equation}
\lambda
=
\frac{1}{2}
\left(
\bar{\sigma}^0{}_{,1}
-
\bar{\sigma}^0
\right)\Omega^{-2}
-
\frac{A\bar{\sigma}^0{}_{,1}}{2}\Omega^{-3}
+
O(\Omega^{-4}).
\tag{7.29}
\end{equation}

\begin{equation}
\begin{aligned}
\nu
={}&
-\frac{1}{2}
\left[
\left(
\bar{\sigma}^0{}_{,1}
-
\bar{\sigma}^0
\right)e^u\nabla S
-
\frac{1}{2}e^uS
\nabla
\left(
\bar{\sigma}^0{}_{,1}
-
\bar{\sigma}^0
\right)
\right]\Omega^{-2}
\\
&+
O(\Omega^{-3}).
\end{aligned}
\tag{7.30}
\end{equation}

\begin{equation}
\begin{aligned}
\alpha
={}&
\frac{1}{2}e^u\bar{\nabla}S\,\Omega^{-2}
-
\frac{1}{2}Ae^u\bar{\nabla}S\,\Omega^{-3}
\\
&+
\frac{1}{4}
\left[
A^2e^u\nabla S
\left(
\frac{S}{2}
+
\frac{1}{2}e^u
\right)
+
e^u(\bar{\nabla}S)\bar{\sigma}^0
\right]\Omega^{-4}
\\
&+
O(\Omega^{-5}).
\end{aligned}
\tag{7.31}
\end{equation}
\begin{equation}
\begin{aligned}
\xi^3
={}&
e^u S\,\Omega^{-2}
-
Ae^u S\,\Omega^{-3}
\\
&+
\frac{1}{2}
\left[
A^2e^u S
\left(
\frac{5}{2}
+
\frac{1}{2}e^{2u}
\right)
-
e^u S\sigma^0
\right]\Omega^{-4}
+
O(\Omega^{-5}).
\end{aligned}
\tag{7.32}
\end{equation}

\begin{equation}
U
=
\frac{1}{2}\Omega^2
-
\frac{1}{8}
\left(
\bar{\psi}_2^0+\psi_2^0
\right)\Omega^{-2}
+
O(\Omega^{-3}).
\tag{7.33}
\end{equation}

\begin{equation}
\chi^3
=
\frac{e^uS}{15}
\left[
\phi_{01}^0
+
\bar{\phi}_{01}^0
+
\psi_1^0
+
\bar{\psi}_1^0
\right]\Omega^{-5}
+
O(\Omega^{-6}\log\Omega).
\tag{7.34}
\end{equation}

\begin{equation}
\begin{aligned}
\psi_4^0
={}&
\left[
-\bar{\sigma}^0
+
\frac{3}{2}\bar{\sigma}^0{}_{,1}
-
\frac{1}{2}\bar{\sigma}^0{}_{,11}
\right]\Omega^{-2}
\\
&+
\left[
-\frac{\bar{\sigma}^0}{2}
-
\frac{3}{4}\bar{\sigma}^0{}_{,1}
+
\frac{1}{2}\bar{\sigma}^0{}_{,11}
\right]\Omega^{-3}
\\
&+
\left[
A^{-2}
\left(
\bar{\sigma}^0
+
\frac{3}{4}\bar{\sigma}^0{}_{,1}
-
\frac{3}{8}\bar{\sigma}^0{}_{,11}
\right)
\right.
\\
&\qquad
+
\frac{A^2}{4}
\left(
-\bar{\sigma}^0
+
\frac{3}{2}\bar{\sigma}^0{}_{,1}
-
\frac{1}{2}\bar{\sigma}^0{}_{,11}
\right)e^{2u}
\\
&\qquad\left.
-
e^{2u}\bar{\nabla}S
\left(
\bar{\sigma}^0{}_{,1}
-
\bar{\sigma}^0
\right)^{3/2}
\nabla
\left[
S
\left(
\bar{\sigma}^0{}_{,1}
-
\bar{\sigma}^0
\right)^{1/2}
\right]
+
\frac{3A}{8}\psi_0^0
\right]\Omega^{-4}
\\
&+
O(\Omega^{-5}).
\end{aligned}
\tag{7.35}
\end{equation}

\begin{equation}
\begin{aligned}
\psi_3
={}&
\left[
e^u
\left(
\bar{\sigma}^0{}_{,1}
-
\bar{\sigma}^0
\right)\nabla S
-
e^uS\nabla
\left(
\bar{\sigma}^0{}_{,1}
-
\bar{\sigma}^0
\right)
\right]\Omega^{-4}
\\
&+
\left[
Ae^u
\left(
-2
\left(
\bar{\sigma}^0{}_{,1}
-
\frac{11}{8}\bar{\sigma}^0
\right)\nabla S
-
S\nabla
\left(
\bar{\sigma}^0{}_{,1}
-
\frac{11}{8}\bar{\sigma}^0
\right)
\right)
+
\frac{1}{2}\bar{\phi}_{01}^0
\right]\Omega^{-5}
\\
&+
O(\Omega^{-6}).
\end{aligned}
\tag{7.36}
\end{equation}


\begin{equation}
\begin{aligned}
\psi_2
={}&
\psi_2^0\Omega^{-6}
+
\Bigg[
-\frac{3}{2}A\psi_2^0
-
e^u S^2\bar{\nabla}
\left(
\frac{\psi_1^0-\phi_{01}^0}{S}
\right)
\\
&\qquad
+
\frac{1}{2}\psi_0^0
\left(
\bar{\sigma}^0{}_{,1}-\bar{\sigma}^0
\right)
+
\frac{3A}{4}\sigma^0\bar{\sigma}^0{}_{,1}
-
\frac{3A}{2}\sigma^0\bar{\sigma}^0
+
16\Lambda^0
\Bigg]\Omega^{-7}
\\
&+
O(\Omega^{-8}\log\Omega).
\end{aligned}
\tag{7.37}
\end{equation}

where:
\begin{equation}
\begin{aligned}
\psi_2^0-\bar{\psi}_2^0
={}&
\Bigg[
e^{2u}
\left(
\frac{S^2}{2}\,
\bar{\nabla}\bar{\nabla}\sigma^0
-
S(\bar{\nabla}S)(\bar{\nabla}\sigma^0)
\right)
\\
&\qquad
+
\sigma^0(\bar{\nabla}S)^2
-
\sigma^0S\bar{\nabla}\bar{\nabla}S
-
\frac{1}{2}\sigma^0\bar{\sigma}^0{}_{,1}
\Bigg]
-\mathrm{c.c.}
\end{aligned}
\tag{7.38}
\end{equation}

\[
\psi_2^0+\bar{\psi}_2^0
\quad\text{is undetermined}.
\]

\[
\psi_1
=
\psi_1^0\Omega^{-7}
+
\psi_1^1\Omega^{-8}\log\Omega
+
\psi_1^2\Omega^{-8}
+
O(\Omega^{-9}\log\Omega).
\]

where
\begin{equation}
\psi_1^0
=
e^u
\left[
S\bar{\nabla}
\left(
\psi_0^0+\frac{15}{4}\sigma^0
\right)
-
\left(
2\psi_0^0+\frac{15}{2}\sigma^0
\right)
\bar{\nabla}S
\right]
-
3\phi_{01}^0.
\tag{7.39}
\end{equation}

\begin{equation}
\psi_1^1
=
-e^u
\left(
S\bar{\nabla}
-
2\bar{\nabla}S
\right)
\left(
5A\psi_0^0
-
\psi_0^1
+
15A\sigma^0
\right)
+
4\phi_{01}^1.
\tag{7.40}
\end{equation}

\[
\psi_1^2
\quad\text{is undetermined}.
\]

By (3.54)
\begin{equation}
\psi_{1,1}^1
=
3\psi_1^1.
\tag{7.41}
\end{equation}


Thus the $u$ derivative of $\psi_1^1$ depends only on itself
and not on the radiation field. It therefore represents a
type of disturbance unconnected with radiation. If it is
zero on one hypersurface, it will remain zero. In this case
it is possible to continue the expansions of all quantities
in negative powers of $\Omega$ without any log terms appearing.

\sectiontoc{The Asymptotic Group}

The metric has the form:
\[
g^{11}=g^{13}=g^{14}=0,
\qquad
g^{12}=1,
\]

\[
g^{22}
=
\Omega^2+O(\Omega^{-2}),
\]

\[
g^{2i}
=
O(\Omega^{-4}),
\]

\[
g^{ij}
=
-2P\bar{P}\,\delta^{ij}\Omega^{-4}
+
2AP\bar{P}\,\delta^{ij}\Omega^{-5}
+
O(\Omega^{-6}).
\]

The asymptotic group is the group of coordinate transformations
that leave the form of the metric and of the boundary conditions
unchanged. It can be derived most simply by considering the
corresponding infinitesimal transformations:

\begin{equation}
x'^{\mu}
=
x^{\mu}
+
\epsilon\,K^{\mu}(x^\nu).
\tag{8.1}
\end{equation}
Then
\begin{equation}
\begin{aligned}
\bar{\delta}g^{\mu\nu}
&=
g'^{\mu\nu}(x)-g^{\mu\nu}(x)
\\
&=
-\epsilon
\left(
g^{\mu\alpha}K^\nu{}_{,\alpha}
+
g^{\nu\alpha}K^\mu{}_{,\alpha}
-
g^{\mu\nu}{}_{,\alpha}K^\alpha
\right).
\end{aligned}
\tag{8.2}
\end{equation}

\begin{equation}
\bar{\delta}\Lambda
=
-\epsilon\,\Lambda_{,\alpha}K^\alpha.
\tag{8.3}
\end{equation}

\begin{equation}
\bar{\delta}\phi_{00}
=
\epsilon R_{1\alpha}K^\alpha{}_{,1}
+
\frac{1}{2}\epsilon R_{11,\alpha}K^\alpha.
\tag{8.4}
\end{equation}

\begin{equation}
\bar{\delta}\phi_{01}
=
\frac{1}{2}\epsilon
\left(
R_{i\alpha}K^\alpha{}_{,1}
+
R_{1\alpha}K^\alpha{}_{,i}
+
R_{1i,\alpha}K^\alpha
\right).
\tag{8.5}
\end{equation}

To obtain the asymptotic group we demand
\begin{equation}
\begin{aligned}
\bar{\delta}g^{11}
&=
\bar{\delta}g^{12}
=
\bar{\delta}g^{13}
=
\bar{\delta}g^{14}
=
0,
\\
\bar{\delta}g^{22}
&=
O(\Omega^{-2}),
\\
\bar{\delta}g^{2i}
&=
O(\Omega^{-4}),
\\
\bar{\delta}g^{ij}
&=
O(\Omega^{-6}),
\\
\bar{\delta}\Lambda
&=
O(\Omega^{-7}),
\\
\bar{\delta}\phi_{00}
&=
O(\Omega^{-9}),
\\
\bar{\delta}\phi_{0i}
&=
O(\Omega^{-7}).
\end{aligned}
\tag{8.6}
\end{equation}

By $\bar{\delta}g^{11}=0$,
\begin{equation}
K^1
=
K^{01}(u,x^i).
\tag{8.7}
\end{equation}

% Continuation of Chapter 3: equations (8.8)--(8.15), Section 9,
% and references. Transcribed from the original thesis scan.

By $\bar{\delta}\Lambda=O(\Omega^{-7})$,
\begin{equation}
K^2=O(\Omega^{-2}).
\tag{8.8}
\end{equation}

\[
\bar{\delta}g^{12}
=0
=K^2{}_{,2}+K^1{}_{,1}
+g^{23}K^1{}_{,3}+g^{24}K^1{}_{,4}.
\]

\[
\therefore\quad K^1{}_{,1}=0,
\]

\begin{equation}
K^1=K^{01}(x^i).
\tag{8.9}
\end{equation}

\[
\bar{\delta}g^{13}
=0
=K^3{}_{,2}+g^{33}K^1{}_{,3}+g^{34}K^1{}_{,4}.
\]

\begin{equation}
K^3
=K^{03}(u,x^i)
-4\bar{P}^{-1}P^2K^1{}_{,3}
+O(\Omega^{-3}).
\tag{8.10}
\end{equation}

\[
\bar{\delta}g^{23}
=O(\Omega^{-4})
=K^3{}_{,1}+O(\Omega^{-4}).
\]

\begin{equation}
K^{03}=K^{03}(x^i).
\tag{8.11}
\end{equation}

\[
\bar{\delta}g^{ij}
=g^{i2}K^j{}_{,2}
+g^{j2}K^i{}_{,2}
+g^{ik}K^j{}_{,k}
+g^{jk}K^i{}_{,k}
-g^{ij}{}_{,k}K^k.
\]

\begin{equation}
K^{03}{}_{,4}=-K^{04}{}_{,3},
\tag{8.12}
\end{equation}

\begin{equation}
K^{03}{}_{,3}=K^{04}{}_{,4},
\tag{8.13}
\end{equation}

\begin{equation}
2S K^{03}{}_{,3}
-S_{,3}K^{03}
-S_{,4}K^{04}
-2S K^1
=0.
\tag{8.14}
\end{equation}

(8.12) and (8.13) imply that $K^{03}$ is an analytic function of
$x^3+i x^4$. This is a consequence of the fact that we have
reduced the leading term of $g^{ij}$ to a conformally flat
form. Thus the only allowed transformations of $x^i$ are the
conformal transformations of the form:

\begin{equation}
x^3+i x^4
=
\frac{a(x^3+i x^4)+b}{c(x^3+i x^4)+d},
\qquad
ad-bc=1.
\tag{8.15}
\end{equation}

When the six parameters $a,b,c,d$ are given $K^1$ is uniquely
determined by (8.14). $K^2$ is also uniquely determined. Thus
the asymptotic group is isomorphic to the conformal group in
two dimensions. Sachs~(7) has shown that this is isomorphic
to the homogeneous Lorentz group. It is also however isomorphic
to the group of motions of a 3-space of constant negative
curvature which is the group of the unperturbed Robertson--
Walker space. Thus the asymptotic group is the same as the
group of the undisturbed space. It is not enlarged by the
presence of radiation. This is interesting because in the
case of gravitational radiation in empty, asymptotically flat
space, it turns out that the asymptotic group contains not
only the 10 dimensional inhomogeneous Lorentz group, the group
of motions of flat space, but also infinite dimensional
``supertranslations''. It has been suggested that these
supertranslations might have some physical significance in
elementary particle physics. The above result would seem
to indicate that this is probably not the case since our
universe is almost certainly not asymptotically flat though
it may be asymptotically Robertson--Walker.

\sectiontoc{What an observer would measure}

The velocity vector $V_m$ of an observer moving with the
dust will be:
\begin{equation}
\begin{aligned}
V_1&=\Omega^{-1}+O(\Omega^{-5}),\\
V_2&=\frac{1}{2}\Omega+O(\Omega^{-3}),\\
V_3&=O(\Omega^{-3}),\\
V_4&=O(\Omega^{-3}).
\end{aligned}
\tag{9.1}
\end{equation}

$q_m$, the projection of the wave vector $l^a$ in the
observer's rest-space (the apparent direction of the wave)
will be:
\[
q_m=\delta_m{}^2-V_mV^1.
\]

\begin{equation}
\begin{aligned}
q_1&=\Omega^{-2}+O(\Omega^{-5}),\\
q_2&=\frac{1}{2}+O(\Omega^{-4}),\\
q_3&=O(\Omega^{-4}),\\
q_4&=O(\Omega^{-4}).
\end{aligned}
\tag{9.2}
\end{equation}

The observer's orthonormal tetrad may be completed by two
space-like unit vectors $s_m$ and $t_m$:
\begin{equation}
\begin{alignedat}{4}
s_1&=O(\Omega^{-4}), &\qquad& t_1&&=O(\Omega^{-4}),\\
s_2&=O(\Omega^{-2}), && t_2&&=O(\Omega^{-2}),\\
s_3&=\frac{1}{\sqrt{2}}, && t_3&&=-\frac{i}{\sqrt{2}},\\
s_4&=\frac{1}{\sqrt{2}}, && t_4&&=\frac{i}{\sqrt{2}}.
\end{alignedat}
\tag{9.3}
\end{equation}

We write
\[
e_a{}^{\alpha}=(V^{\alpha},q^{\alpha},s^{\alpha},t^{\alpha}).
\]

By measuring the relative accelerations of neighbouring dust
particles, the observer may determine the `electric'
components of the gravitational wave:
\[
E_{ab}=-C_{apbq}V^pV^q.
\]

In the observer's tetrad this has components
\begin{equation}
\begin{aligned}
E={}&
\begin{pmatrix}
0&0&0\\
0&1&0\\
0&0&-1
\end{pmatrix}O(\Omega^{-4})
+
\begin{pmatrix}
0&0&0\\
0&0&1\\
0&1&0
\end{pmatrix}O(\Omega^{-4})
\\[1ex]
&+
\begin{pmatrix}
0&1&1\\
1&0&0\\
1&0&0
\end{pmatrix}O(\Omega^{-5})
+
\begin{pmatrix}
0&1&-1\\
1&0&0\\
-1&0&0
\end{pmatrix}O(\Omega^{-5})
\\[1ex]
&+
\begin{pmatrix}
1&0&0\\
0&-\frac{1}{2}&0\\
0&0&-\frac{1}{2}
\end{pmatrix}O(\Omega^{-6}).
\end{aligned}
\tag{9.4}
\end{equation}

This should be compared to the behaviour for asymptotically
flat space for which
\begin{equation}
\begin{aligned}
E={}&
\begin{pmatrix}
0&0&0\\
0&1&0\\
0&0&-1
\end{pmatrix}O(r^{-1})
+
\begin{pmatrix}
0&0&0\\
0&0&1\\
0&1&0
\end{pmatrix}O(r^{-1})
\\[1ex]
&+
\begin{pmatrix}
0&1&1\\
1&0&0\\
1&0&0
\end{pmatrix}O(r^{-2})
+
\begin{pmatrix}
0&1&-1\\
1&0&0\\
-1&0&0
\end{pmatrix}O(r^{-2})
\\[1ex]
&+
\begin{pmatrix}
1&0&0\\
0&-\frac{1}{2}&0\\
0&0&-\frac{1}{2}
\end{pmatrix}O(r^{-3}).
\end{aligned}
\tag{9.5}
\end{equation}

\sectiontoc{References}

\begin{enumerate}
\item H. Bondi, M. G. J. van der Burg and A. W. K. Metzner,
  \emph{Proc. Roy. Soc. A} \textbf{269}, 21 (1962).

\item R. K. Sachs,
  \emph{Proc. Roy. Soc. A} \textbf{270}, 103 (1962).

\item E. T. Newman and R. Penrose,
  \emph{J. Math. Phys.} \textbf{3}, 566 (1962).

\item E. T. Newman and T. W. J. Unti,
  \emph{J. Math. Phys.} \textbf{3}, 891 (1962).

\item D. Norman, Thesis, London University (1963).

\item R. Penrose,
  \emph{Proc. Roy. Soc. A} \textbf{284}, 159 (1965).

\item R. K. Sachs,
  \emph{Phys. Rev.} \textbf{128}, 2851 (1962).
\end{enumerate}
